\documentclass[11pt,a4paper]{article}
\usepackage[utf8]{inputenc}
\usepackage[T1]{fontenc}
\usepackage[margin=2.3cm]{geometry}
\usepackage{booktabs,longtable,array,amsmath,enumitem}
\usepackage{xurl}
\usepackage[hidelinks]{hyperref}
\setlist{itemsep=1pt,topsep=3pt}
\newcolumntype{L}[1]{>{\raggedright\arraybackslash}p{#1}}
\newcommand{\est}{\textsc{est}}
\newcommand{\ver}{\textsc{ver}}
\title{HL-LHC electroweak quirks, study v2\\[2pt]\large Daily report, day 2}
\author{Private work (simulation). Public information only.}
\date{18 November 2026, 18:55 (day 2 of the approved calendar)}
\begin{document}
\maketitle

\noindent\textbf{Tags.} \ver{} = measured or read in this study, with the campaign or file in the repository.
\est{} = estimate. All reach and background numbers are \est{} on v1 placeholder efficiencies. No v2 signal
efficiency has been computed; the sealed v1 signal counts of the reproduction test stay sealed.

\section{Summary}
\begin{itemize}
\item \textbf{Waiting for you:} nothing. The two small choices of yesterday (trigger turn-on uncertainty,
      thickness-law diagnostic) come with recommendations at decision point 1 on 21 November.
\item You approved amendment E4 at 12:44 as recommended. With it the SM family of the reproduction test passes.
      The frozen FAIL stays on record, labelled ``after the result''.
\item Pilot part 1 is at 457 of 460 Z($\nu\nu$) jobs (minimum bias complete); part 2 at 664 of 762 jobs, all
      class blocks done. No pilot job has failed.
\item The new Opus orchestrator session started this morning. One reproduction campaign (C1 rates) had failed
      unseen on 17 November; it was rerun and passes.
\end{itemize}

\section{What finished today}
\begin{enumerate}
\item \textbf{Per-stream validation of the single-interaction pilot.} 8 of 8 streams pass under the approved
      reading (D9: events with hits before their production time are flagged, not dropped) (\ver). One line for
      you, as the advisor asked: pilot-wide there are \textbf{318 flagged events} (524 hits) in
      $1.72\times10^{7}$, against the ``about 170'' in the notice on which you decided D9. They appear in every
      nuclear-interaction class, not only in the one seen before. D9 stands: it is a per-event rule with no
      ceiling. The preregistration of v2 will state which variant (with or without these events) is binding;
      the recommendation is \emph{with}, the other as a cross-check.
\item \textbf{C1 rates} (rerun after a missing generator setting): 5 of 5 pass (\ver). The job script now
      reports failed jobs instead of finishing silently.
\item \textbf{SM family (Z($\nu\nu$) candidates per event).} Failed as frozen at rung R04: $+2.6\,\%$,
      $z=+3.64$. A regression ticket and an independent advisor found the cause in the error, not the physics:
      $1/\sqrt{N}$ treats the $\approx47$ candidates of one event as independent. With event-level errors
      (amendment E4): $z=+0.35/+0.57/+0.50$ at R04--R06, block $\chi^2=1.53/3$; with v1's error set to zero,
      at most $+0.82$ (\ver). The test is weak: a real 15\,\% shift at R04--R05 would be seen only about half the
      time.
\item \textbf{Plane-fit accuracy, second half (E2 restricted scan).} Fails: 46 of 50 and 46 of 50 against the
      required 50 (\ver). The first half (100 of 100) stands alone, as E2 foresaw.
\item \textbf{C2-b passes; C2-c passes but has no power} (flagged as such) (\ver).
\item \textbf{Sample stitching on $10^{5}$ events:} ratio $1.0001\pm0.0252$ (\ver). The 12\,\% excess seen
      on 20{,}000 events was a fluctuation; $f_{200}$ is not biased.
\item \textbf{Pile-up family (C$'_{\mathrm{PU}}$):} the same error correction was committed \emph{before} its
      numbers are read (commit \texttt{54610b6}). A partial file computed earlier with the old error is
      superseded and was not used.
\end{enumerate}

\section{Decisions for you}
None today. Coming at decision point 1 (21 November), each with a recommendation checked by an advisor:
the trigger turn-on uncertainty and the thickness-law diagnostic.

\section{Running now}
\begin{center}
\begin{tabular}{L{3.6cm}L{10.6cm}}
\toprule
Pilot, part 1 & Z($\nu\nu$) 457 of 460 jobs on schrodinger, two lanes; ends about 20:00 tonight (\est) \\
Pilot, part 2 & 664 of 762 jobs. The hard-scatter block (98 jobs) starts after the last Z($\nu\nu$) file;
                through about 20 November (\est). Its validation job is queued behind it \\
v1 reproduction & Pile-up $\rho=1$: 173 of 200 (ends about 19 November 04:30); pile-up pairs: 49 of 75
                (overnight). Then a new executor reads the pile-up family and all 41 comparisons \\
Generator comparison & Waits for a runner change, made after the campaigns above are complete \\
\bottomrule
\end{tabular}
\end{center}
No agent is running tonight.

\section{Budget and storage}
\begin{center}
\begin{tabular}{L{7.6cm}L{6.4cm}}
\toprule
CPU used (ledger, 2{,}768 jobs) & 383 CPU-h of 4{,}609 approved (\ver) \\
\quad of which reproduction / theory / pilot channel / pilot & 195 / 39 / 13 / 136 CPU-h \\
Projected total & 4{,}663--4{,}844 CPU-h (\est), plus 22--28 for D2; threshold 5{,}992 \\
Storage, this study & 236 GB of 1.5 TB \\
Shared volume free & 3.7 TB of 4.5 TB \\
\bottomrule
\end{tabular}
\end{center}
The projection is redone with measured costs at the checkpoint of day 5.

\section{Problems}
\begin{itemize}
\item \textbf{Unseen failure.} The 12 C1-rate jobs failed on 17 November (generator data path not set) and the
      queue did not report it. Fixed and rerun; the queue now reports failures.
\item \textbf{Timestamps.} Three log entries this morning carried times written without reading the clock; a
      correction entry gives the true times (10:42--10:48).
\item \textbf{Open before the full reproduction read:} the coherence test of the pile-up and closure families
      still treats nested rungs as independent; it is fixed before those numbers are read.
\item \textbf{Failed jobs since day 0:} 36 failed, 3 timed out, 1 killed, all in development and reproduction
      campaigns (the 12 above included); none in the pilot.
\end{itemize}

\section{Reach: where it stands (\est)}
Unchanged: Gate A expected a reach below 243~GeV, near the existing bounds of about 180~GeV, and a pause at
day 15 ($B_{50}$ expected $3\times10^{4}$--$3\times10^{5}$ against the cap of 300). Nothing measured today tests
it. The stitching result removes one downward risk on the signal side. The first measured background number comes
with the preliminary projection of day 7.

\section{Plan for 19--20 November}
\begin{itemize}
\item Pilot part 1 completes; the hard-scatter block of part 2 runs, then its validation job.
\item Reproduction campaigns complete; a new executor fixes the family coherence, reads the pile-up family and all
      41 comparisons.
\item Then the sealed comparison (advisor first), and the decision-point-1 package for 21 November.
\end{itemize}

\section{Your next interventions}
\begin{center}
\begin{tabular}{L{4.2cm}L{10.0cm}}
\toprule
21 November (day 5) & Decision point 1: reproduction result; two small choices with recommendations \\
23 November (day 7) & Preliminary background projection (to read); possible Z($\nu\nu$) top-up decision \\
1 December (day 15) & Binding background projection; a pause is expected \\
\bottomrule
\end{tabular}
\end{center}

\section*{Files}
\begin{itemize}
\item \path{analyses/quirks_v2/phase2_exploration/outputs/B1_STREAM_VALIDATION_hedda_2026-11-18.md}
\item \path{analyses/quirks_v2/phase2_exploration/REGRESSION_TICKET_SM.md}
\item \path{analyses/quirks_v2/phase2_exploration/outputs/SM_FIX_lisbet_2026-11-18.md}
\item \path{analyses/quirks_v2/phase2_exploration/outputs/C2_VERDICTS_ignatius_2026-11-18.md}
\item \path{analyses/quirks_v2/phase2_exploration/outputs/STITCH_100k_ignatius_2026-11-18.md}
\item \path{analyses/quirks_v2/COMMITMENTS.md} (amendment E4)
\item \path{analyses/quirks_v2/ORCHESTRATOR_STATE.md}
\end{itemize}
\end{document}
